% GENERATED FILE. Edit canonical lessons, then run npm run generate:books.
% canonical-source-hash: fnv1a64:bb4661202beffedc
% canonical-lessons: LA-C335-auriga, LA-C335-quadrigae, LA-C335-meta, LA-C335-factio, LA-C335-circensis

\chapter{Lūdī circēnsēs}
\label{ch:la-ludi-circenses}
\hlchaptermodality{335}

\begin{chapteropening}
\textbf{By the end of this chapter:} \emph{I can describe a chariot race: the charioteer, the four-horse chariot, the turning post, the racing team and of the circus.}

Complete the last lesson of chapter 335: I can describe a chariot race: the charioteer, the four-horse chariot, the turning post, the racing team and of the circus.
\end{chapteropening}

\section[aurīga, aurīgae]{aurīga, aurīgae — a charioteer}

\label{lesson:LA-C335-auriga}



\begin{quote}
Before the new one: say the Latin for to vow, then the Latin for a vow.
\end{quote}

\subsection*{You'll want to know: aurīga, aurīgae}
\textbf{aurīga, aurīgae} — \textquotedblleft{}a charioteer\textquotedblright{}.

\begin{quote}
Aurīga equōs celeriter agēbat.
\end{quote}

\begin{quote}
\emph{The charioteer was driving the horses fast.}
\end{quote}

Masculine, 1st declension, though it ends in \emph{-a}, like \emph{nauta}. Famous charioteers grew very rich.

\begin{cousinweb}
\textbf{aurīga} is explained by Lewis \& Short as built on an old word for the reins and \emph{agō}, \textquotedblleft{}to drive\textquotedblright{}: he who handles the reins. The details are not settled.
\end{cousinweb}

\begin{grammarlens}[title={What you've built}]
One more word for telling, explaining and describing at B1.
\end{grammarlens}

\subsection*{Your turn}
\begin{itemize}
\raggedright
  \item \emph{Say it:} \emph{aurīga, aurīgae}
  \item \emph{Say it:} \emph{aurīga, aurīgae}, once more
  \item \emph{Say it:} a sentence of your own with \emph{aurīga, aurīgae}
  \item \emph{Recall:} say the Latin for to vow, then the Latin for a vow, then say \emph{aurīga, aurīgae} again
\end{itemize}

\begin{tcolorbox}[breakable,colback=teal!4,colframe=teal!35!black,title={Before you move on}]
What does aurīga, aurīgae mean? (\textquotedblleft{}A charioteer\textquotedblright{}.) Say it once more.
\end{tcolorbox}
\section[quadrīgae, quadrīgārum]{quadrīgae, quadrīgārum — a four-horse chariot}

\label{lesson:LA-C335-quadrigae}



\begin{quote}
Before the new one: say the Latin for a vow, then the Latin for a charioteer.
\end{quote}

\subsection*{You'll want to know: quadrīgae, quadrīgārum}
\textbf{quadrīgae, quadrīgārum} — \textquotedblleft{}a four-horse chariot\textquotedblright{}.

\begin{quote}
Quadrīgae per circum volābant.
\end{quote}

\begin{quote}
\emph{The four-horse chariots flew round the circus.}
\end{quote}

Feminine, plural: four horses under one yoke.

\begin{cousinweb}
\textbf{quadrīgae} is a shortened \emph{quadriiugae}, from \emph{quattuor}, \textquotedblleft{}four\textquotedblright{}, and \emph{iugum}, \textquotedblleft{}yoke\textquotedblright{} (Lewis \& Short): four together.
\end{cousinweb}

\begin{grammarlens}[title={What you've built}]
One more word for telling, explaining and describing at B1.
\end{grammarlens}

\subsection*{Your turn}
\begin{itemize}
\raggedright
  \item \emph{Say it:} \emph{quadrīgae, quadrīgārum}
  \item \emph{Say it:} \emph{quadrīgae, quadrīgārum}, once more
  \item \emph{Say it:} a sentence of your own with \emph{quadrīgae, quadrīgārum}
  \item \emph{Recall:} say the Latin for a vow, then the Latin for a charioteer, then say \emph{quadrīgae, quadrīgārum} again
\end{itemize}

\begin{tcolorbox}[breakable,colback=teal!4,colframe=teal!35!black,title={Before you move on}]
What does quadrīgae, quadrīgārum mean? (\textquotedblleft{}A four-horse chariot\textquotedblright{}.) Say it once more.
\end{tcolorbox}
\section[mēta, mētae]{mēta, mētae — a turning post; a goal, a limit}

\label{lesson:LA-C335-meta}



\begin{quote}
Before the new one: say the Latin for a charioteer, then the Latin for a four-horse chariot.
\end{quote}

\subsection*{You'll want to know: mēta, mētae}
\textbf{mēta, mētae} — \textquotedblleft{}a turning post; a goal, a limit\textquotedblright{}.

\begin{quote}
Aurīga prope mētam cecidit.
\end{quote}

\begin{quote}
\emph{The charioteer fell near the turning post.}
\end{quote}

Feminine. The cone-shaped posts at each end of the track, where crashes happened. Also any limit or goal.

\begin{cousinweb}
\textbf{mēta} is inherited from Proto-Italic \emph{*mētā}, \textquotedblleft{}boundary\textquotedblright{} (EtymDB): a marker at the limit.
\end{cousinweb}

\begin{grammarlens}[title={What you've built}]
One more word for telling, explaining and describing at B1.
\end{grammarlens}

\subsection*{Your turn}
\begin{itemize}
\raggedright
  \item \emph{Say it:} \emph{mēta, mētae}
  \item \emph{Say it:} \emph{mēta, mētae}, once more
  \item \emph{Say it:} a sentence of your own with \emph{mēta, mētae}
  \item \emph{Recall:} say the Latin for a charioteer, then the Latin for a four-horse chariot, then say \emph{mēta, mētae} again
\end{itemize}

\begin{tcolorbox}[breakable,colback=teal!4,colframe=teal!35!black,title={Before you move on}]
What does mēta, mētae mean? (\textquotedblleft{}A turning post; a goal, a limit\textquotedblright{}.) Say it once more.
\end{tcolorbox}
\section[factiō, factiōnis]{factiō, factiōnis — a racing team; a party, a faction}

\label{lesson:LA-C335-factio}



\begin{quote}
Before the new one: say the Latin for a four-horse chariot, then the Latin for a turning post.
\end{quote}

\subsection*{You'll want to know: factiō, factiōnis}
\textbf{factiō, factiōnis} — \textquotedblleft{}a racing team; a party, a faction\textquotedblright{}.

\begin{quote}
Puer factiōnem viridem amābat.
\end{quote}

\begin{quote}
\emph{The boy loved the Green team.}
\end{quote}

Feminine, from \emph{faciō}. The racing teams wore colours: white, red, green and blue. English \emph{faction}.

\begin{cousinweb}
\textbf{factiō} is built on \emph{faciō}, \textquotedblleft{}to make, to do\textquotedblright{} (EtymDB): a doing together, and so a group that acts together. English \emph{faction} comes from it.
\end{cousinweb}

\begin{grammarlens}[title={What you've built}]
One more word for telling, explaining and describing at B1.
\end{grammarlens}

\subsection*{Your turn}
\begin{itemize}
\raggedright
  \item \emph{Say it:} \emph{factiō, factiōnis}
  \item \emph{Say it:} \emph{factiō, factiōnis}, once more
  \item \emph{Say it:} a sentence of your own with \emph{factiō, factiōnis}
  \item \emph{Recall:} say the Latin for a four-horse chariot, then the Latin for a turning post, then say \emph{factiō, factiōnis} again
\end{itemize}

\begin{tcolorbox}[breakable,colback=teal!4,colframe=teal!35!black,title={Before you move on}]
What does factiō, factiōnis mean? (\textquotedblleft{}A racing team; a party, a faction\textquotedblright{}.) Say it once more.
\end{tcolorbox}
\section[circēnsis, circēnse]{circēnsis, circēnse — of the circus}

\label{lesson:LA-C335-circensis}



\begin{quote}
Before the new one: say the Latin for a turning post, then the Latin for a racing team.
\end{quote}

\subsection*{You'll want to know: circēnsis, circēnse}
\textbf{circēnsis, circēnse} — \textquotedblleft{}of the circus\textquotedblright{}.

\begin{quote}
Lūdī circēnsēs per quīnque diēs fīēbant.
\end{quote}

\begin{quote}
\emph{The races in the circus went on for five days.}
\end{quote}

Two endings. \emph{Lūdī circēnsēs}: the races. A satirist said the people wanted only bread and races.

\begin{cousinweb}
\textbf{circēnsis} is built on \emph{circus}, \textquotedblleft{}ring, racecourse\textquotedblright{} (EtymDB), with the ending \emph{-ēnsis}, \textquotedblleft{}belonging to\textquotedblright{}.
\end{cousinweb}

\begin{grammarlens}[title={What you've built}]
One more word for telling, explaining and describing at B1.
\end{grammarlens}

\subsection*{Your turn}
\begin{itemize}
\raggedright
  \item \emph{Say it:} \emph{circēnsis, circēnse}
  \item \emph{Say it:} \emph{circēnsis, circēnse}, once more
  \item \emph{Say it:} a sentence of your own with \emph{circēnsis, circēnse}
  \item \emph{Recall:} say the Latin for a turning post, then the Latin for a racing team, then say \emph{circēnsis, circēnse} again
\end{itemize}

\begin{tcolorbox}[breakable,colback=teal!4,colframe=teal!35!black,title={Before you move on}]
What does circēnsis, circēnse mean? (\textquotedblleft{}Of the circus\textquotedblright{}.) Say it once more.
\end{tcolorbox}
